% Restyle of fig-sealed-composition-crossover to the tikz-diagram-craft v2
% standard (measured-quantity plot; the plot itself is unchanged).
%
% Reader question: when does the advanced composition bound start to pay?
% Claim: below k=32 releases (at the chapter's per-release epsilon=0.1) the
% basic bound is smaller and should be quoted; past it, the advanced bound is.
% Evidence: closed forms k*epsilon and sqrt(2 k ln(1/delta')) epsilon +
% k epsilon (e^epsilon-1), website-v2/public/whitepaper/sealed-harbor.tex,
% Sec.~sealed-conservation, "Numbers by hand" [verified].
% Must distinguish: the two curves; the crossover point; which bound is
% smaller on each side of it.
% Grammar chosen: measured-quantity plot (pgfplots), the worked point (k=32)
% marked on both curves -- a schematic curve pair would lose the fact that
% the crossover is a specific, computed k.
% Grammar rejected: a table of the two bounds at k=32 alone would drop the
% crossover structure -- that a contract's right bound depends on how many
% releases it already expects.
% Five-second test: which curve is lower before k=32, and which after.
\begin{figure}[H]
\centering
\pdfigurehue{pdcobalt}%
\begin{tikzpicture}[pd figure]
  \begin{axis}[
    width=9.6cm,height=5.6cm,
    axis lines=left,
    title={which bound to quote, as the engagement grows longer},
    title style={font=\pdfiglabelbold,text=pdink,yshift=-2pt},
    xmin=0,xmax=64,ymin=0,ymax=7,
    xlabel={releases $k$ (per-release $\varepsilon=0.1$)},
    ylabel={spent-budget bound},
    xlabel style={font=\pdfiglabel,text=pdinkmuted},
    ylabel style={font=\pdfiglabel,text=pdinkmuted},
    tick label style={font=\pdfiglabel,text=pdinkmuted},
    axis line style={pd hairline},
    tick style={pd tick},
    xtick={0,16,32,48,64},
    ytick={0,2,4,6},
    clip=false,
    legend style={at={(0.03,0.97)},anchor=north west,draw=none,fill=none,
      font=\pdfiglabel,text=pdink,row sep=1pt,inner sep=2pt},
    legend cell align=left,
  ]
    \addplot[domain=0:64,samples=120,pd rule] {0.1*x};
    \addlegendentry{basic: $k\varepsilon$}
    \addplot[domain=0:64,samples=120,pd focus rule] {sqrt(2*x*ln(1/0.000001))*0.1 + x*0.1*(exp(0.1)-1)};
    \addlegendentry{advanced: $\sqrt{k}$-scaled}
    \addplot[only marks,mark=*,mark size=1.4pt,draw=pdink,fill=pdink,forget plot] coordinates {(32,3.20)};
    \addplot[only marks,mark=*,mark size=1.4pt,draw=pdfocus,fill=pdfocus,forget plot] coordinates {(32,3.3101)};
    \draw[pd guide] (axis cs:32,0) -- (axis cs:32,3.3101);
    \node[pd tag,align=left,text width=2.4cm,anchor=north west] at (axis cs:34.5,2.7) {$k=32$:\\basic $3.20$\\advanced $3.31$};
  \end{axis}
\end{tikzpicture}
\caption{The basic and advanced $\varepsilon$-composition bounds on cumulative spend as functions of the number of releases $k$, at the per-release $\varepsilon=0.1$ the chapter uses for this comparison. At $k=32$, the chapter's reported comparison, basic composition ($3.20$) is still smaller than the advanced bound ($3.31$): below the crossover a contract should quote the exact, $\delta$-free basic bound, and only past it the square-root-scaled advanced bound -- the same rule at $\varepsilon=0.05$ gives $k=128$: basic $6.40$ versus advanced $3.30$, where the advanced bound has come to pay.}
\label{fig:sealed-composition-crossover}
\end{figure}
